\documentclass[12pt,a4paper]{ctexart}

\usepackage[margin=2.2cm]{geometry}
\usepackage{amsmath,amssymb,mathtools}
\usepackage{array,booktabs,longtable}
\usepackage[dvipsnames]{xcolor}
\usepackage{hyperref}
\usepackage{listings}
\usepackage{tikz}
\usetikzlibrary{arrows.meta,calc,fit,positioning}

\hypersetup{
  colorlinks=true,
  linkcolor=blue!60!black,
  urlcolor=blue!60!black
}

\lstdefinestyle{py}{
  language=Python,
  basicstyle=\ttfamily\small,
  keywordstyle=\color{blue!70!black},
  commentstyle=\color{ForestGreen!70!black},
  stringstyle=\color{BrickRed},
  showstringspaces=false,
  breaklines=true,
  frame=single,
  rulecolor=\color{black!20},
  backgroundcolor=\color{black!2}
}

\newcommand{\filepath}[1]{\texttt{#1}}
\newcommand{\code}[1]{\texttt{#1}}
\newcommand{\vect}[1]{\boldsymbol{#1}}

\title{第 3 章注意力机制学习教程\\\large 基于 \filepath{ch03/01\_main-chapter-code/ch03.ipynb}}
\author{}
\date{\today}

\begin{document}

\maketitle
\tableofcontents
\newpage
\bigskip

\section{这份教程解决什么问题}

原始 notebook \filepath{ch03/01\_main-chapter-code/ch03.ipynb} 的目标不是把你变成
``会背 attention 公式的人''，而是把你带到一个更实用的位置：

\begin{quote}
你能从零解释 self-attention，能把它写成 PyTorch 模块，能看懂 causal mask，
也能理解为什么 GPT 里的注意力最后会变成 multi-head attention。
\end{quote}

这份中文教程不逐段翻译原文，而是按学习路径重新组织：

\begin{enumerate}
  \item 为什么需要 attention；
  \item 最小版 self-attention 到底怎么算；
  \item 为什么真实模型要引入 \(W_q\)、\(W_k\)、\(W_v\)；
  \item 为什么语言模型必须加 causal mask；
  \item multi-head attention 究竟比单头多了什么；
  \item 学完这一章后，下一步能做什么输出。
\end{enumerate}

\section{先看全局主线}

这一章的核心不是很多零散知识点，而是一条连续主线：

\begin{center}
\begin{tikzpicture}[
  node distance=1.1cm,
  stage/.style={draw, rounded corners, fill=blue!6, minimum width=4.2cm, minimum height=0.9cm, align=center},
  arr/.style={-Latex, thick}
]
\node[stage] (s1) {长序列建模困难};
\node[stage, below=of s1] (s2) {先做无参数的 self-attention\\建立直觉};
\node[stage, below=of s2] (s3) {引入 \(W_q, W_k, W_v\)\\得到可训练注意力};
\node[stage, below=of s3] (s4) {加入 causal mask 和 dropout\\变成语言模型可用版本};
\node[stage, below=of s4] (s5) {把单头扩展成 multi-head\\接入后续 GPT};

\draw[arr] (s1) -- (s2);
\draw[arr] (s2) -- (s3);
\draw[arr] (s3) -- (s4);
\draw[arr] (s4) -- (s5);
\end{tikzpicture}
\end{center}

如果你在阅读时丢了方向，就回到上面这条链：这一章一直在回答
``为什么这样设计 attention 模块''，而不只是教你把几行 PyTorch 代码敲出来。

\section{为什么长序列建模会逼出 attention}

\subsection{问题的本质}

在早期序列模型里，常见做法是把整句信息压进一个固定大小的隐状态，再靠这个隐状态继续生成输出。
这会遇到两个实际问题：

\begin{enumerate}
  \item 输入一长，前面信息容易被后面过程稀释；
  \item 不同输出位置需要依赖输入中的不同部分，但固定向量很难对这种选择性依赖建模。
\end{enumerate}

attention 的想法更直接：

\begin{quote}
当前 token 需要什么信息，就去全序列里按相关性取，而不是指望一个固定摘要把所有信息都无损保存。
\end{quote}

\subsection{学习时应抓住的转变}

从这里开始，你要把自己从``压缩整句''的思维切到``按需检索上下文''的思维。
后面的所有公式，本质上都只是在实现下面这句话：

\begin{quote}
给定当前位置，计算它和其他位置的相关性，再用这些相关性对其他位置的信息加权求和。
\end{quote}

\section{第一步：无可训练参数的 self-attention}

\subsection{这一步的作用}

原 notebook 先故意讲一个\textbf{不真实但很适合建立直觉}的版本：
直接拿输入向量自己和自己做点积，不引入任何可学习参数。

这样做的好处是，你可以先理解 attention 的计算链，而不被模型参数分散注意力。

\subsection{示例输入}

notebook 用一句非常短的 toy sentence 做例子：

\begin{center}
\texttt{Your journey starts with one step}
\end{center}

每个 token 先被表示成一个 3 维向量。示例矩阵如下：

\begin{center}
\begin{tabular}{c c c}
\toprule
token & 记号 & toy embedding \\
\midrule
Your & \(x^{(1)}\) & \([0.43, 0.15, 0.89]\) \\
journey & \(x^{(2)}\) & \([0.55, 0.87, 0.66]\) \\
starts & \(x^{(3)}\) & \([0.57, 0.85, 0.64]\) \\
with & \(x^{(4)}\) & \([0.22, 0.58, 0.33]\) \\
one & \(x^{(5)}\) & \([0.77, 0.25, 0.10]\) \\
step & \(x^{(6)}\) & \([0.05, 0.80, 0.55]\) \\
\bottomrule
\end{tabular}
\end{center}

\subsection{三步计算链}

假设我们现在只关心第二个 token \(x^{(2)}\)。
这一节的目标是计算它的上下文向量 \(z^{(2)}\)。

\begin{enumerate}
  \item 用 \(x^{(2)}\) 作为查询，与所有输入向量做点积，得到 attention scores；
  \item 对 scores 做 softmax，得到总和为 1 的 attention weights；
  \item 用 attention weights 对所有输入向量做加权求和，得到 context vector。
\end{enumerate}

\begin{center}
\begin{tikzpicture}[
  box/.style={draw, rounded corners, minimum width=2.1cm, minimum height=0.9cm, align=center},
  small/.style={draw, rounded corners, minimum width=1.2cm, minimum height=0.7cm, align=center, fill=blue!5},
  arr/.style={-Latex, thick}
]
\node[box, fill=yellow!15] (query) at (0, 0) {查询\\\(x^{(2)}\)};
\node[small] (x1) at (-2.8, -1.8) {\(x^{(1)}\)};
\node[small, fill=yellow!15] (x2) at (-1.6, -1.8) {\(x^{(2)}\)};
\node[small] (x3) at (-0.4, -1.8) {\(x^{(3)}\)};
\node[small] (x4) at (0.8, -1.8) {\(x^{(4)}\)};
\node[small] (x5) at (2.0, -1.8) {\(x^{(5)}\)};
\node[small] (x6) at (3.2, -1.8) {\(x^{(6)}\)};

\node[box, fill=green!10] (score) at (0, -3.5) {点积\\scores};
\node[box, fill=orange!12] (softmax) at (0, -5.0) {softmax\\weights};
\node[box, fill=red!8] (ctx) at (0, -6.6) {加权求和\\\(z^{(2)}\)};

\draw[arr] (query) -- (score);
\foreach \n in {x1,x2,x3,x4,x5,x6} {
  \draw[arr] (\n) -- (score);
}
\draw[arr] (score) -- (softmax);
\draw[arr] (softmax) -- (ctx);
\foreach \n in {x1,x2,x3,x4,x5,x6} {
  \draw[arr] (\n) |- (ctx);
}
\end{tikzpicture}
\end{center}

\subsection{这一步真正要学会的东西}

这一节最重要的不是背住符号，而是把下面三件事说清楚：

\begin{enumerate}
  \item \textbf{attention score} 是相关性的粗分数，还没有归一化；
  \item \textbf{attention weight} 是 softmax 后的权重，总和为 1；
  \item \textbf{context vector} 不是复制原 token，而是融合了其他 token 信息的新表示。
\end{enumerate}

也就是说，attention 的输出不是``把 query 挑出来''，而是
``围绕 query 重新组织一遍整个上下文''。

\subsection{为什么原 notebook 单独讲 naive softmax}

notebook 先写一个最朴素的 softmax：

\begin{lstlisting}[style=py]
def softmax_naive(x):
    return torch.exp(x) / torch.exp(x).sum(dim=0)
\end{lstlisting}

然后再切回 \code{torch.softmax}。这不是废话，而是在提醒你：

\begin{enumerate}
  \item softmax 本质上只是把分数转成概率型权重；
  \item 真实实现里，数值稳定性和库实现细节非常重要。
\end{enumerate}

\section{第二步：把 attention 推广到整句}

前一节只算了一个 \(z^{(2)}\)。但语言模型不会只更新一个位置，它要同时为每个 token 计算上下文向量。

这时就从``一个 query 对全体 key''，自然推广到``每个位置都做一次 query''。
原 notebook 做了一个很重要的思维跳转：

\begin{quote}
原本看起来像双层循环的计算，其实可以整体改写成矩阵乘法。
\end{quote}

于是：
\[
\text{attn\_scores} = X X^\top, \qquad
\text{attn\_weights} = \operatorname{softmax}(\text{attn\_scores}), \qquad
Z = \text{attn\_weights} X
\]

学习时你要特别注意这一点，因为后面所有高效实现都依赖这个视角：
attention 看起来是在两两比较 token，真正落地时则是\textbf{张量批量运算}。

\section{第三步：引入可训练的 Q、K、V 投影}

\subsection{为什么前面的版本不够}

如果直接用输入向量自己互相点积，模型其实没有学会``应该用什么标准来判断相关性''。
它只是把现成 embedding 空间里的几何关系拿来用。

真实模型需要更强的表达能力，于是引入三组可训练参数：
\[
q^{(i)} = x^{(i)} W_q, \qquad
k^{(i)} = x^{(i)} W_k, \qquad
v^{(i)} = x^{(i)} W_v
\]

\subsection{这三个量各自干什么}

\begin{enumerate}
  \item Query：当前这个位置正在\textbf{找什么}；
  \item Key：每个位置手里拿着一张\textbf{可被匹配的标签}；
  \item Value：真正会被聚合到输出里的\textbf{信息载体}。
\end{enumerate}

一个常见误区是把它们当成三份随机复制品。
更准确的理解是：

\begin{quote}
同一个输入 \(x\) 被投影进了三种不同职责的空间：
匹配空间（query/key）和传递信息空间（value）。
\end{quote}

\subsection{为什么要除以根号 d\_k}

notebook 在这里引入 scaled dot-product attention：
\[
\operatorname{softmax}\!\left(\frac{QK^\top}{\sqrt{d_k}}\right)V
\]

要点不是记住这个形式，而是理解原因：

\begin{enumerate}
  \item 维度大时，点积数值会变大；
  \item 直接送进 softmax 容易让分布过尖；
  \item 过尖会让梯度不稳定，训练效果变差；
  \item 所以用 \(\sqrt{d_k}\) 做尺度校正。
\end{enumerate}

\subsection{原 notebook 里的两个类}

\filepath{ch03.ipynb} 先写 \code{SelfAttention\_v1}，再写 \code{SelfAttention\_v2}。
它们的核心逻辑一致，差别主要在实现风格：

\begin{enumerate}
  \item \code{SelfAttention\_v1} 直接用 \code{nn.Parameter} 自己管理权重矩阵；
  \item \code{SelfAttention\_v2} 改用 \code{nn.Linear}，更接近真实工程写法。
\end{enumerate}

其中 \code{SelfAttention\_v2} 的核心非常短：

\begin{lstlisting}[style=py]
class SelfAttention_v2(nn.Module):
    def __init__(self, d_in, d_out, qkv_bias=False):
        super().__init__()
        self.W_query = nn.Linear(d_in, d_out, bias=qkv_bias)
        self.W_key   = nn.Linear(d_in, d_out, bias=qkv_bias)
        self.W_value = nn.Linear(d_in, d_out, bias=qkv_bias)

    def forward(self, x):
        keys = self.W_key(x)
        queries = self.W_query(x)
        values = self.W_value(x)

        attn_scores = queries @ keys.T
        attn_weights = torch.softmax(
            attn_scores / keys.shape[-1]**0.5, dim=-1
        )
        context_vec = attn_weights @ values
        return context_vec
\end{lstlisting}

学习时应该盯住 \code{forward}，不是盯住类名。
如果你能解释这一段每一行在干什么，这一节就算学到了点子上。

\section{第四步：为什么语言模型必须使用 causal attention}

\subsection{问题出在哪里}

前面的 self-attention 默认一个 token 可以看到全序列所有位置。
这对编码任务通常没问题，但对自回归语言模型是致命的：

\begin{quote}
如果当前位置在训练时看到了未来 token，模型就等于作弊。
\end{quote}

因此，语言模型里的注意力必须满足：

\begin{center}
位置 \(t\) 只能看见 \(1\) 到 \(t\) 的 token，不能看见 \(t+1\) 之后的 token。
\end{center}

\subsection{mask 的几何形状}

这就是为什么会出现一个上三角 mask。下面的示意图里，绿色区域表示允许关注，
红色区域表示未来位置，必须屏蔽：

\begin{center}
\begin{tikzpicture}[scale=0.9]
\foreach \r in {0,1,2,3,4} {
  \foreach \c in {0,1,2,3,4} {
    \pgfmathtruncatemacro{\isfuture}{\c>\r ? 1 : 0}
    \ifnum\isfuture=1
      \fill[red!20] (\c,-\r) rectangle ++(1,-1);
    \else
      \fill[green!14] (\c,-\r) rectangle ++(1,-1);
    \fi
    \draw[black!40] (\c,-\r) rectangle ++(1,-1);
  }
}
\node at (2.5,0.7) {列：被看的位置};
\node[rotate=90] at (-0.7,-2.5) {行：当前 query 位置};
\node[anchor=west] at (5.4,-1.2) {\color{green!50!black}\rule{0.8cm}{0.25cm} 可见};
\node[anchor=west] at (5.4,-2.0) {\color{red!60!black}\rule{0.8cm}{0.25cm} 屏蔽未来};
\end{tikzpicture}
\end{center}

\subsection{为什么要在 softmax 前屏蔽}

原 notebook 先演示了一个直观但不够优雅的方案：
先 softmax，再把未来位置乘零，然后重新归一化。

紧接着它说明了更标准的做法：

\begin{quote}
在 softmax 前，把未来位置的 score 直接填成 \(-\infty\)。
\end{quote}

这样 softmax 后这些位置自然变成 0，而且每一行仍然保持合法概率分布。

\subsection{dropout 在这里扮演什么角色}

dropout 不是 causal mask 的替代品，它解决的是另一个问题：训练时的过拟合。
notebook 选择在 attention weights 上施加 dropout，这是常见做法。

你可以把这一步理解成：

\begin{quote}
模型在训练时不要过度依赖某几个固定连接，而要学会更稳健地分配注意力。
\end{quote}

\subsection{批处理版本的意义}

到 \code{CausalAttention} 这一类时，notebook 不再只处理单句，而是把输入 shape 扩成
\((b, \text{num\_tokens}, d)\)。

这一点非常关键，因为它标志着代码开始从``教学算例''进入``可接入真实训练循环''的阶段。

\section{第五步：从单头到 multi-head attention}

\subsection{为什么不能只用一个头}

单头 attention 的问题不在于它不能用，而在于它太容易把所有关系都压进一个注意力模式里。
multi-head 的核心好处是：

\begin{quote}
让模型在多个子空间里并行地学习不同类型的依赖关系。
\end{quote}

比如一个头偏向局部语法，一个头偏向长距离依赖，另一个头偏向指代或结构信息。

\subsection{直观结构图}

\begin{center}
\begin{tikzpicture}[
  box/.style={draw, rounded corners, minimum width=2.0cm, minimum height=0.9cm, align=center},
  arr/.style={-Latex, thick}
]
\node[box, fill=blue!8] (x) at (0,0) {输入 \(X\)};
\node[box, fill=green!10] (qkv) at (0,-1.5) {线性投影\\得到 \(Q,K,V\)};
\node[box, fill=orange!12] (split) at (0,-3.1) {按 head 切分};

\node[box, fill=yellow!15] (h1) at (-3.3,-4.8) {Head 1\\attention};
\node[box, fill=yellow!15] (h2) at (0,-4.8) {Head 2\\attention};
\node[box, fill=yellow!15] (h3) at (3.3,-4.8) {Head \(h\)\\attention};

\node[box, fill=red!10] (concat) at (0,-6.5) {拼接 heads};
\node[box, fill=purple!10] (proj) at (0,-8.0) {输出投影};

\draw[arr] (x) -- (qkv);
\draw[arr] (qkv) -- (split);
\draw[arr] (split) -- (h1);
\draw[arr] (split) -- (h2);
\draw[arr] (split) -- (h3);
\draw[arr] (h1) -- (concat);
\draw[arr] (h2) -- (concat);
\draw[arr] (h3) -- (concat);
\draw[arr] (concat) -- (proj);
\end{tikzpicture}
\end{center}

\subsection{两种实现思路}

原 notebook 先给了 \code{MultiHeadAttentionWrapper}，它的想法很直观：
直接堆多个单头 \code{CausalAttention}。

然后才给更标准的 \code{MultiHeadAttention}：

\begin{enumerate}
  \item 先一次性投影到 \(d_{\text{out}}\)；
  \item 再 reshape 成 \((b, \text{num\_tokens}, \text{num\_heads}, \text{head\_dim})\)；
  \item 再 transpose 到 \((b, \text{num\_heads}, \text{num\_tokens}, \text{head\_dim})\)；
  \item 每个 head 独立算 attention；
  \item 拼回去，再做输出投影。
\end{enumerate}

\subsection{这部分真正难的不是公式，而是 shape}

很多人第一次学 multi-head attention 觉得抽象，不是因为公式难，而是因为张量 shape 一直在变。
所以学习这一节时，最有价值的训练是把每一步 shape 写出来。

例如 notebook 里的核心流向是：
\[
\begin{aligned}
X &: (b, n, d_{\text{in}}) \\
Q, K, V &: (b, n, d_{\text{out}}) \\
\text{reshape 后} &: (b, n, h, d_h) \\
\text{transpose 后} &: (b, h, n, d_h) \\
\text{scores} &: (b, h, n, n) \\
\text{context} &: (b, h, n, d_h) \\
\text{拼回输出} &: (b, n, d_{\text{out}})
\end{aligned}
\]

只要你把这组 shape 走顺了，multi-head attention 就不会再神秘。

\section{第六步：如何把这一章学成可输出的能力}

\subsection{学完后最应该做的三种输出}

\paragraph{输出一：口头解释}
你应该能不用看 notebook，独立解释下面这些问题：

\begin{enumerate}
  \item attention score 和 attention weight 有什么区别；
  \item 为什么要有 \(Q\)、\(K\)、\(V\) 三组投影；
  \item 为什么要用 \(\sqrt{d_k}\) 缩放；
  \item causal mask 为什么必须在 softmax 前介入；
  \item multi-head 相比单头到底多了什么表达能力。
\end{enumerate}

\paragraph{输出二：最小实现}
你应该能在空白文件里手写一个最小版注意力模块：

\begin{lstlisting}[style=py]
class TinyAttention(nn.Module):
    def __init__(self, d_in, d_out):
        super().__init__()
        self.Wq = nn.Linear(d_in, d_out, bias=False)
        self.Wk = nn.Linear(d_in, d_out, bias=False)
        self.Wv = nn.Linear(d_in, d_out, bias=False)

    def forward(self, x):
        q = self.Wq(x)
        k = self.Wk(x)
        v = self.Wv(x)
        scores = q @ k.transpose(-2, -1)
        weights = torch.softmax(scores / k.shape[-1]**0.5, dim=-1)
        return weights @ v
\end{lstlisting}

\paragraph{输出三：接入后续 GPT}
这一章的 \code{MultiHeadAttention} 不是孤立玩具。
在仓库后续章节里，它会进入 GPT 模型主体，例如：

\begin{enumerate}
  \item \filepath{ch04/01\_main-chapter-code/gpt.py}
  \item 多个 \filepath{previous\_chapters.py} 文件
\end{enumerate}

也就是说，这一章学完以后，你已经具备了理解 GPT block 中注意力子模块的基础。

\section{建议的学习节奏}

\subsection{第一遍：只抓主线}

第一遍读 \filepath{ch03.ipynb} 时，只做一件事：
把整章分成五个阶段记住。

\begin{enumerate}
  \item 无参数 self-attention；
  \item 矩阵化推广；
  \item 可训练 \(Q/K/V\)；
  \item causal mask + dropout；
  \item multi-head attention。
\end{enumerate}

\subsection{第二遍：盯 shape 和 forward}

第二遍不要再看大段说明，直接盯几个关键实现：

\begin{enumerate}
  \item \code{SelfAttention\_v2}
  \item \code{CausalAttention}
  \item \code{MultiHeadAttention}
\end{enumerate}

你的目标不是再看懂一次，而是把每个 \code{forward} 里的张量流向写在纸上。

\subsection{第三遍：做一个小闭环}

建议你做这三个动作：

\begin{enumerate}
  \item 自己手写最小 attention；
  \item 打印每一步 shape；
  \item 去看第 4 章的 GPT 实现，确认 attention 模块是怎么嵌进去的。
\end{enumerate}

\section{本章最短总结}

如果把 \filepath{ch03.ipynb} 压成几句话，它讲的其实就是：

\begin{quote}
attention 的本质是：先比较相关性，再按相关性取信息。\\
self-attention 让每个 token 都能看全局。\\
引入 \(Q/K/V\) 后，这种比较和取信息的方式变成了可学习的。\\
加入 causal mask 后，attention 才能用于自回归语言模型。\\
加入 multi-head 后，模型能在多个子空间中并行学习不同关系。\\
最后，这个模块会直接进入后续 GPT 的实现。
\end{quote}

\bigskip
\noindent
\textbf{建议的下一步：}
读完本教程后，回到原 notebook，重点看下面三个实现：

\begin{enumerate}
  \item \code{SelfAttention\_v2}
  \item \code{CausalAttention}
  \item \code{MultiHeadAttention}
\end{enumerate}

然后自己写一个最小版实现。这样你学到的就不再只是概念，而是可以输出的能力。

\end{document}
